% ==================== CHƯƠNG 11 ====================
\newpage
\chap{Hạn chế và hướng phát triển}\label{sec:limits}

\sect{Hạn chế}
\begin{enumerate}
    \item \textbf{Dữ liệu}: chỉ có giai đoạn 2011--2012, quy mô vừa phải (17.379 dòng sự kiện), mức tổng toàn hệ thống; không có vị trí trạm, số xe sẵn sàng, lịch sử điều chuyển xe. Vì vậy DSS chỉ khuyến nghị mức độ chuẩn bị, chưa giải được bài toán điều phối giữa các trạm.
    \item \textbf{Chất lượng dữ liệu}: 165 giờ thiếu bản ghi (76 ngày không đủ 24 giờ), 22 giờ độ ẩm bằng 0, 12,5\% số giờ có tốc độ gió bằng 0, mức thời tiết 4 chỉ có 3 giờ; mã mùa của nguồn lệch so với mùa theo cảm nhận (``Spring'' thực chất là mùa đông). Các vấn đề này đã được ghi nhận và gắn cờ, nhưng không sửa được từ phía nguồn.
    \item \textbf{Nạp dữ liệu}: ETL nạp lại toàn bộ (phù hợp dữ liệu tĩnh), chưa có cơ chế nạp tăng dần, chưa có lịch chạy tự động.
    \item \textbf{Mô hình DSS}: dựa trên trung bình lịch sử theo nhóm; chưa xét đồng thời nhiều biến liên tục (nhiệt độ, độ ẩm, gió). Kiểm định ngoài mẫu (mục~\ref{sec:dss-backtest}) cho thấy mô hình tĩnh bỏ sót nhiều giờ cao điểm khi nhu cầu tăng (recall mức Cao 58,5\%), và hiệu chỉnh mức theo sai số gần nhất khắc phục phần lớn (recall 90,0\%). Hạn chế còn lại: kiểm định dùng thời tiết thực tế của năm 2012 nên kết quả vận hành sẽ thấp hơn, và mới chỉ có một năm kiểm thử.
    \item \textbf{Hệ thống}: chưa có xác thực người dùng, chưa triển khai lên máy chủ; tương quan nhiệt độ, độ ẩm chưa kiểm soát đầy đủ biến gây nhiễu.
\end{enumerate}

\sect{Hướng phát triển}
\begin{enumerate}
    \item \textbf{Mở rộng dữ liệu}: dùng dữ liệu chuyến đi chi tiết mà Capital Bikeshare công bố (có trạm đi, trạm đến, thời lượng), bổ sung chiều \textbf{Trạm} (áp dụng SCD loại 2 cho thay đổi sức chứa, vị trí) và bảng sự kiện \textbf{tồn xe theo trạm}; khi đó DSS có thể tính số xe cần điều chuyển giữa các trạm (bài toán tối ưu, ví dụ quy hoạch tuyến tính).
    \item \textbf{Nạp tăng dần}: ETL theo lô hằng ngày, dùng khóa nghiệp vụ \code{(date\_key, hour\_key)} để upsert; phân vùng bảng sự kiện theo năm khi dữ liệu lớn; lập lịch chạy tự động.
    \item \textbf{Mô hình dự báo}: thay mô hình tra cứu bằng mô hình hồi quy hoặc học máy (ví dụ gradient boosting) nhận đầu vào giờ, loại ngày, dự báo thời tiết, nhiệt độ, xu hướng; phân loại mức dựa trên khoảng dự báo thay vì chỉ trung bình. Khung kiểm định ngoài mẫu đã có (mục~\ref{sec:dss-backtest}: huấn luyện 2011, kiểm tra 2012, cùng các chỉ số độ chính xác, recall mức Cao, tỷ lệ thiếu/dư chuẩn bị, MAE), nên có thể dùng lại để so sánh mô hình mới với mô hình hiện tại trên cùng thước đo.
    \item \textbf{Tích hợp dự báo thời tiết} qua API để tự động lập kế hoạch cho ngày mai và gửi cảnh báo (email, tin nhắn) cho người vận hành.
    \item \textbf{Hệ thống}: thêm xác thực, phân quyền; đóng gói Docker và triển khai; bổ sung bộ đệm (cache) cho các truy vấn tổng hợp hoặc materialized view khi dữ liệu lớn.
\end{enumerate}
